\section{Model Architecture Evolution}
\label{sec:expansion}

The model grows through three types of architectural expansion, applied at
specific stages.  Each expansion is validated by measuring cosine similarity
of output logits between the pre- and post-expansion models.

\subsection{Depth Expansion (Layer Cloning)}

\textbf{Applied at:} Stages 2 (6$\to$9 layers) and 4 (9$\to$12 layers).

Selected layers are deep-copied and appended to the model.  Symmetry-breaking
Gaussian noise ($\sigma = 0.01$) is added to all weight matrices (not biases
or norms) in the cloned layers.

\textbf{Why symmetry breaking is mandatory.}
Without noise, cloned layers produce identical gradients and learn identical
functions, providing no capacity benefit.  The noise ensures gradient
differentiation from the first update.

\textbf{Clone selection.}
Layers from the upper half of the network are cloned (e.g., layers [3,4,5] for
6$\to$9).  Upper layers capture more task-specific features, so cloning them
provides useful initialization for learning new domain representations.

\textbf{This is not function-preserving.}
Appending active Transformer blocks changes the model output.  Validation shows
cosine similarity $\approx 0.92$ after depth expansion---indicating significant
but controlled drift.

\subsection{FFN Widening}

\textbf{Applied at:} Stage 6 ($d_\text{ff}$: 2048$\to$3584).

The FFN is widened with function-preserving initialization:
\begin{itemize}[nosep]
  \item $W_\text{gate}$ and $W_\text{up}$: old rows copied, new rows initialized
        with small random values ($\sigma = 0.02/\sqrt{2 \cdot n_\text{layers}}$).
  \item $W_\text{down}$: old columns copied, new columns set to \textbf{zero}.
\end{itemize}
The zero-initialized $W_\text{down}$ columns ensure the new neurons are
\emph{dormant} at initialization---they contribute nothing to the output.
This makes the expansion \textbf{exactly function-preserving}:
$f(x; \theta_\text{expanded}) = f(x; \theta_\text{original})$ for all $x$.

\subsection{Context Length Extension}

\textbf{Applied at:} Stages 2 (256$\to$384), 4 (384$\to$512), and 6
(512$\to$768).

For learnable positional embeddings, the embedding matrix
$\mathbf{P} \in \mathbb{R}^{L_\text{old} \times d}$ is extended to
$\mathbf{P}' \in \mathbb{R}^{L_\text{new} \times d}$ via 1D linear
interpolation (see Section~\ref{sec:posenc}).  This is approximately
function-preserving for positions within the original range.

\subsection{Why Changes Are Staged}

All three expansion types could be applied simultaneously, but staging them
provides several benefits:
\begin{enumerate}[nosep]
  \item \textbf{Controlled drift.}  Each expansion introduces a bounded amount
        of output perturbation.  The training loop can recover from one change
        before the next is applied.
  \item \textbf{Diagnostic clarity.}  If training diverges or forgetting spikes,
        the cause can be attributed to the specific expansion.
  \item \textbf{Curriculum alignment.}  Each expansion is paired with a new
        domain that benefits from the added capacity (e.g., longer context for
        longer stories, wider FFN for richer vocabulary).
\end{enumerate}

\begin{table}[H]
\centering\small
\caption{Architecture at each stage.}
\begin{tabular}{clcccc}
\toprule
\textbf{Stage} & \textbf{Expansion} & \textbf{Layers} & $d_\textbf{ff}$ &
\textbf{ctx} & \textbf{Params} \\
\midrule
0 & (baseline)                   & 6  & 2048 & 256 & 45.8M \\
2 & +3 layers, ctx 384           & 9  & 2048 & 384 & 58.4M \\
3 & (data-only)                  & 9  & 2048 & 384 & 58.4M \\
4 & +3 layers, ctx 512           & 12 & 2048 & 512 & 71.1M \\
5 & (data-only)                  & 12 & 2048 & 512 & 71.1M \\
6 & FFN 3584, ctx 768            & 12 & 3584 & 768 & 99.5M \\
\bottomrule
\end{tabular}
\label{tab:stages}
\end{table}
